% SGA 3，第 VI_B 篇，第一轮英文重建稿的中文翻译。
% 范围：引理10.3至引理10.8及其间全部证明、推论10.5、命题10.6和编者注93--94。
% 依据：Polo--Gille Exp6B-13oct24.pdf，本组件局部页70--71
% （合卷阅读版页411--412；再版页眉400--401；原版页边印刷页387--390）。
% 引理10.9始于局部页71，不在本组件内。
% 编者注95属于被排除的引理10.9之后的10.10，因此不在本组件内。
% OCR及Reinhold英文Markdown仅用作对照见证。
% 第二轮以忠实的原生LaTeX重建替换了图(D)的源页裁切图；替换已完成。

\medskip
\phantomsection\label{lem:sga3-VIB-10.3}%
\noindent\textbf{引理 10.3.---}
\emph{设 \(G\) 是有限呈示的 \(S\)-群。
则存在指标 \(j\in I\)、有限呈示的 \(S_j\)-群 \(G_j\)，以及 \(S\)-群同构}
\[
  G\xrightarrow{\ \sim\ }G_j\times_{S_j}S.
\]

考虑到 EGA~0\(_{\mathrm{III}}\)，8.2.5和8.2.6中对 \(S\)-群结构的解释，%
\footnote{\label{note:sga3-VIB-93}编者注：即用 \(S\)-态射的交换图来解释。}
断言 \hyperref[lem:sga3-VIB-10.2]{10.2} 和
\hyperref[lem:sga3-VIB-10.3]{10.3} 容易由
\hyperref[par:sga3-VIB-10.0]{10.0} 和
\hyperref[par:sga3-VIB-10.1]{10.1} 推出。

\medskip
\phantomsection\label{lem:sga3-VIB-10.4}%
\noindent\textbf{引理 10.4.---}
\emph{设 \(u:G\to H\) 是有限呈示的 \(S\)-群之间的 \(S\)-群态射。
由 \hyperref[lem:sga3-VIB-10.3]{10.3} 和
\hyperref[lem:sga3-VIB-10.2]{10.2}，\(u\) 由有限呈示的
\(S_j\)-群之间的态射 \(u_j:G_j\to H_j\) 经基变换得到。
则 \(u\) 是中心单态射\textup{（分别地，正规单态射）}，当且仅当存在
\(i\geq j\)，使得}
\phantomsection\label{form:sga3-VIB-10.4-base-change}%
\[
  u_i=u_j\times_{S_j}S_i
\]
\emph{具有相同性质。}

由 \hyperref[rem:sga3-VIB-6.7]{6.7} 中对中心或正规群单态射的刻画，
这立即由 \hyperref[par:sga3-VIB-10.0]{10.0} 和
\hyperref[par:sga3-VIB-10.1]{10.1} 得到。

\medskip
\phantomsection\label{cor:sga3-VIB-10.5}%
\noindent\textbf{推论 10.5.---}
\emph{设 \(G_j\) 是有限呈示的 \(S_j\)-群。
则 \(G_j\times_{S_j}S\) 交换，当且仅当存在 \(i\geq j\)，使得
\(G_j\times_{S_j}S_i\) 交换。}

事实上，说一个 \(S\)-群交换，等价于说将其视为自身的群子概形时，它是中心的。

\medskip
\phantomsection\label{prop:sga3-VIB-10.6}%
\noindent\textbf{命题 10.6.---}
\emph{设 \(G_j\) 是有限呈示的 \(S_j\)-群，\(G'_j\) 是
\(G_j\) 的一个群子概形，在 \(S_j\) 上平坦且有限呈示。则商}
\phantomsection\label{form:sga3-VIB-10.6-quotient-over-S}%
\[
  (G_j\times_{S_j}S)/(G'_j\times_{S_j}S)
\]
\emph{对于 \textup{(fpqc)} 拓扑可表，当且仅当存在 \(i\geq j\)，使}
\phantomsection\label{form:sga3-VIB-10.6-quotient-over-Si}%
\[
  (G_j\times_{S_j}S_i)/(G'_j\times_{S_j}S_i)
\]
\emph{可表。}

这是下面这个更一般的引理的推论。

\medskip
\phantomsection\label{lem:sga3-VIB-10.7}%
\noindent\textbf{引理 10.7.---}
\emph{设 \(X_j\) 是有限呈示的 \(S_j\)-概形，\(R_j\) 是
\(X_j\) 上的一个平坦且有限呈示的等价关系。}%
\footnote{\label{note:sga3-VIB-94}编者注：即复合态射
\[
  R_j\lhook\joinrel\longrightarrow X_j\times_{S_j}X_j
  \xrightarrow{\ \operatorname{pr}_1\ }X_j
\]
是平坦且有限呈示的。}
\emph{则对于拓扑 \(\mathcal T=\textup{(fppf)}\) 或
\(\textup{(fpqc)}\) 的商层}
\[
  (X_j\times_{S_j}S)/(R_j\times_{S_j}S)
\]
\emph{可表，当且仅当存在 \(i\geq j\)，使对于拓扑 \(\mathcal T\) 的商层}
\[
  (X_j\times_{S_j}S_i)/(R_j\times_{S_j}S_i)
\]
\emph{可表。}

由 \hyperref[par:sga3-VIB-10.0]{10.0} 中回顾的 EGA~IV\(_3\)，\footnote{规范源注
（SGA-SRC-000861）：现行再版和官方原版均印作IV\(_2\)，但所引四节均属于EGA~IV\(_3\)。}
8.8.2、8.8.3、8.10.5和11.2.6中的论断，本引理可由下一个结果推出。

\medskip
\phantomsection\label{lem:sga3-VIB-10.8}%
\noindent\textbf{引理 10.8.---}
\emph{设 \(\mathcal T\) 是 \textup{(fppf)} 或 \textup{(fpqc)} 拓扑。
设 \(X\) 是有限呈示\textup{（分别地，局部有限呈示）}的 \(S\)-概形，
\(R\) 是 \(X\) 上由单态射}
\[
  v:R\longrightarrow X\times_S X
\]
\emph{定义的等价关系，使 \(\operatorname{pr}_1\circ v\) 平坦且有限呈示
\textup{（分别地，平坦且局部有限呈示）}。则下列条件等价：}
\begin{enumerate}
\item[\textup{(i)}]\phantomsection\label{concl:sga3-VIB-10.8-i}
  \emph{对于拓扑 \(\mathcal T\) 的商层 \(X/R\) 可表。}

\item[\textup{(ii)}]\phantomsection\label{concl:sga3-VIB-10.8-ii}
  \emph{存在有限呈示\textup{（分别地，局部有限呈示）}的 \(S\)-概形
  \(Y\)，以及忠实平坦态射 \(p:X\to Y\)，使下图}
\end{enumerate}
\begin{center}
  \phantomsection\label{diag:sga3-VIB-10.8-D}%
  \begin{tabular}{c@{\qquad}c}
    \textup{(D)} &
    \begin{tikzcd}[row sep=large,column sep=large]
      R
        \arrow[r,"{\operatorname{pr}_1\circ v}"]
        \arrow[d,"{\operatorname{pr}_2\circ v}"'] &
      X \arrow[d,"p"] \\
      X \arrow[r,"p"'] &
      Y
    \end{tikzcd}
  \end{tabular}
\end{center}
\emph{是笛卡尔图。}

首先注意，由IV，\SGARefLink{sga3:iv:definition:3:3:2}{3.3.2}和4.4.3，
拓扑 \(\mathcal T\) 的层 \(X/R\) 由 \(Y\) 表示，当且仅当图
\hyperref[diag:sga3-VIB-10.8-D]{\textup{(D)}} 是笛卡尔图，
且 \(p\) 是拓扑 \(\mathcal T\) 的覆盖。

下面证明 \hyperref[concl:sga3-VIB-10.8-i]{\textup{(i)}} 蕴含
\hyperref[concl:sga3-VIB-10.8-ii]{\textup{(ii)}}。
假设 \hyperref[concl:sga3-VIB-10.8-i]{\textup{(i)}} 蕴含图
\hyperref[diag:sga3-VIB-10.8-D]%
{\SGARefLink{diag:sga3-VIB-10.8-D}{\textup{(D)}}} 为笛卡尔图。
因此 \(\operatorname{pr}_1\circ v\) 由 \(p\) 沿 \(p\) 作基变换得到，
且 \(p\) 是 \textup{(fpqc)} 拓扑的覆盖。
从而，由 \textup{(fpqc)} 下降（EGA~IV\(_2\)，2.7.1），
既然 \(\operatorname{pr}_1\circ v\) 忠实平坦且\textup{（局部）}有限呈示，
\(p\) 也如此。于是 EGA~IV\(_4\)，17.7.5 表明，由于 \(X\)
在 \(S\) 上\textup{（局部）}有限呈示，\(Y\) 也如此。

下面证明 \hyperref[concl:sga3-VIB-10.8-ii]{\textup{(ii)}} 蕴含
\hyperref[concl:sga3-VIB-10.8-i]{\textup{(i)}}。
只须证明 \(p\) 是 \textup{(fppf)} 拓扑的覆盖。
但根据假设，\(p\) 忠实平坦；又因 \(X\) 和 \(Y\) 在 \(S\) 上
局部有限呈示，\(p\) 也局部有限呈示
\textup{（EGA~IV\(_1\)，1.4.3\textup{(v)}）}。
